\section{Introduction}

Transient hardware faults remain an important threat to computing-system reliability. Continued device scaling, higher integration density, and lower operating voltages may increase hardware susceptibility to soft errors \cite{mukherjee2005soft,borkar2006designing,baumann2005radiation,frustaci2015sram}, while physical drift in emerging memory devices may also induce data bit flips \cite{yilmaz2017drift,junsangsri2014new}. These faults may manifest as silent data corruption (SDC), producing incorrect results without triggering observable failures. This concern is particularly important for large language models (LLMs). Built primarily on the Transformer architecture \cite{vaswani2017attention}, LLMs have demonstrated strong capabilities across language understanding, generation, and reasoning tasks \cite{brown2020language,zhao2023survey}. During autoregressive inference, model layers are executed repeatedly, allowing perturbations in parameters, activations, or intermediate results to propagate across decoding steps. Moreover, increasing model sizes are driving LLM inference toward multi-GPU deployment, where faults may also occur and propagate across devices.

Existing neural-network reliability studies have developed various fault injection tools for convolutional neural networks (CNNs) and conventional deep neural networks (DNNs) \cite{mahmoud2021optimizing,wan2024mulberry,jiao2017assessment,liu2022special,xu2021reliability,reagen2018ares,mahmoud2020pytorchfi,huang2024mrfi}. However, these tools provide limited native support for important characteristics of modern LLM inference, including distributed parallelism, quantized model representations, and iterative decoding. In contrast, vLLM incorporates parallelism, quantization, and system optimizations within a common inference architecture \cite{kwon2023efficient}, making it a suitable foundation for fault injection under practical LLM inference configurations.

Two additional limitations prevent traditional approaches from scaling effectively to distributed LLM inference. First, existing model-level fault models commonly determine the number of computation faults from the number of elements in an operation's output tensor. Output size, however, does not reflect the underlying computational exposure: operations with the same output dimensions may require substantially different numbers of multiply--accumulate operations. Furthermore, communication-buffer faults are rarely covered, leaving an important fault source in multi-GPU inference insufficiently represented. Second, a common injection path transfers selected GPU tensor values to the CPU for bit-level modification and then copies them back to the GPU. This implementation is simple and can be acceptable for relatively small CNN and DNN workloads with limited forward execution and sparse fault injection. For LLMs, however, large tensors, iterative decoding, and increasing fault counts repeatedly introduce host--device transfers and synchronization, causing the injection overhead to grow rapidly. Direct GPU modification avoids this path but requires careful handling of data-type representations, concurrent updates, and asynchronous CUDA execution.

To address these challenges, we present vLLM-FI, a fault-exposure-aware runtime fault injection and reliability evaluation framework integrated with vLLM. vLLM-FI extends model-level fault injection to cover memory, computation, and communication faults. Its computation fault model accounts for computational exposure beyond output tensor size, while its communication fault model supports injection into software-visible communication buffers. To reduce injection overhead, vLLM-FI uses customized CUDA operators to perform in-place bit flips on GPU-resident tensors. Because vLLM executes models through distributed GPU workers, vLLM-FI employs a configuration-driven architecture to coordinate fault settings across processes and resolve them to the intended injection locations. The framework supports tensor parallelism, pipeline parallelism, and multiple full-precision and quantized model representations. It also integrates with lm-evaluation-harness, connecting configurable fault injection experiments with standardized downstream-task evaluation.

The main contributions of this paper are summarized as follows:

\begin{itemize}
    \item We develop fault-exposure-aware models for distributed LLM inference, covering memory, computation, and communication faults. The models derive compute-fault counts from runtime multiply--accumulate workloads rather than output size alone, and communication-fault counts from buffer sizes in tensor-parallel collectives and pipeline-stage transfers.

    \item We design a vLLM-native fault injection architecture that supports distributed and quantized LLM inference. Customized CUDA operators perform efficient in-place bit flips on GPU-resident tensors, while configuration-driven targeting coordinates fault injection across distributed execution processes.

    \item We integrate vLLM-FI with lm-evaluation-harness to provide an end-to-end workflow from fault configuration and injection to standardized reliability evaluation. We validate the framework and analyze LLM reliability across fault models, downstream tasks, quantization strategies, and model-parallel configurations. vLLM-FI is released as open source to support reproducible research.
\end{itemize}
